\section{Structure de la Solution et Implémentation Itérative}

Un \textit{matching engine} suit la logique du \textit{price-time priority}: 
\begin{itemize}
    \item les ordres d'achat les plus élevés sont priorisés par rapport à ceux plus faibles
    \item les ordres de vente les plus bas sont priorisés par rapport à ceux plus élevés
\end{itemize}

Lorsque le prix de vente le plus faible correspond au prix d'achat le plus élevé (par comparaison d'infériorité), un \textit{match} a lieu et la transaction prend place, au prix de l'ordre déjà présent dans le carnet (l'ordre \textit{résident}). Pour deux \textit{matchs} équivalents en termes de prix, la priorité est donnée selon une condition temporelle (premier arrivé, premier servi). Un \textit{matching engine} dans l'industrie est responsable de plusieurs tâches : vérifier que les prix d'échange ne sont pas complètement décorrélés du marché, s'assurer d'appliquer l'appétence au risque de l'entreprise par rapport aux quantités et au prix global des actifs achetés, etc. Pour ce mémoire, nous nous contenterons d'appliquer le \textit{price-time priority} sur des données ingérées depuis un générateur artificiel et nous excluerons toutes les fonctionnalités annexes d'un \textit{matching engine}.

\subsection{Conception générale du \textit{matching engine}}

\subsubsection{Contexte métier et technique}

Notre \textit{matching engine} est un \textit{Proof of Concept} qui simule le cœur d'une plateforme d'échange électronique. Les acteurs du système sont de deux natures : en amont, un \textbf{générateur d'ordres synthétiques} qui représente les participants de marché (\textit{market makers}, traders algorithmiques) en émettant des ordres d'achat et de vente ; en aval, un \textbf{système de persistance} qui consomme les transactions exécutées, jouant le rôle d'un système de reporting réglementaire.

L'enjeu qui justifie ce mémoire est le suivant : les marchés financiers modernes exigent une réactivité que les architectures logicielles classiques, telles qu'une JVM standard, ne garantissent pas nativement. Chaque microseconde de latence représente un risque financier quantifiable, notamment en cas de mouvement de marché brutal. Le \textit{matching engine} doit donc assurer la fiabilité et le déterminisme des transactions sans intervention humaine, et ce, en continu (le moteur s'exécute sans interruption planifiée). Un jeu de données de test sera utilisé pour garantir la consistance des résultats entre les différentes versions du moteur. L'enjeu central de ce mémoire reste l'optimisation de la latence et du déterminisme du programme.

\subsubsection{Exigences fonctionnelles du \textit{matching engine}}

Le \textit{matching engine} permet strictement de :

\begin{itemize}
    \item générer des ordres d'achat et de vente de \textit{commodities} en temps réel (avec prix et quantité)
    \item annuler des ordres émis
    \item actualiser le prix d'un ordre émis
    \item autoriser des \textit{matchs} partiels (par exemple, si un acteur veut vendre 10 unités et qu'un autre souhaite en acheter 5, le \textit{match} a lieu sur 5 unités, et l'ordre de vente reste actif pour les 5 unités restantes)
    \item faire le \textit{matching} entre les ordres selon le \textit{price-time priority}
    \item persister les transactions effectuées
\end{itemize}

Tout ordre qui ne trouve pas immédiatement de contrepartie n'est pas rejoué : il reste dans le carnet d'ordres (à l'état \textit{resting}) en attendant un \textit{match} futur ou une annulation. Il s'agit du comportement nominal d'un carnet d'ordres.

Le \textit{matching engine} ne gère pas :
\begin{itemize}
    \item les types d'ordres autres que le \textit{limit order} (\textit{market order}, \textit{stop order}, \textit{iceberg order} sont exclus)
    \item le multi-instruments : le système opère sur un instrument unique (mono-instrument). L'extension à plusieurs instruments impliquerait un carnet d'ordres et potentiellement un thread de \textit{matching} dédiés par instrument.
    \item la fermeture des marchés (il tourne en continu)
    \item le contrôle de cohérence entre le prix d'un ordre et le marché, ni la gestion du risque (appétence au risque, exposition globale)
    \item la sécurité des transactions et l'authentification des acteurs
    \item la tolérance aux pannes et la réplication
\end{itemize}

\subsubsection{Exigences non-fonctionnelles du \textit{matching engine}}\label{exigences-nf}
Les exigences suivantes constituent la \textbf{cible visée pour la version finale optimisée} du moteur. La \textit{baseline} (Jalon 0), dépourvue de toute optimisation, est attendue de dépasser significativement ces seuils — notamment à cause des pauses \textit{Stop-The-World} du \textit{Garbage Collector} — et c'est précisément cet écart, mesuré jalon après jalon, qui constitue le fil conducteur de l'implémentation.

\begin{itemize}
    \item \textbf{Déterminisme} : un même jeu de données d'entrée doit produire exactement les mêmes \textit{matchs}, dans le même ordre.
    \item \textbf{Latence} : latence P99.995 de bout en bout de l'ordre de la cinquantaine de microsecondes, et une latence maximale inférieure à la milliseconde, pour la version finale.
    \item \textbf{Débit} : la baseline est testée à une charge de référence de 10 000 opérations par seconde. La version finale optimisée vise un débit cible de l'ordre du million d'opérations par seconde — un objectif ambitieux mais ancré dans la littérature industrielle (LMAX revendique 6 millions d'ordres par seconde sur un seul thread en régime pleinement optimisé).
    \item \textbf{Bornage des prix} : les prix sont bornés à un intervalle fixe, par exemple $[0, 100\,000]$ exprimés en centimes (soit $0$ à $1\,000{,}00$ unités monétaires). Ce bornage est une hypothèse de conception qui justifiera, dans les jalons avancés, l'encodage des prix en entiers (arithmétique à virgule fixe) plutôt qu'en \texttt{double}.
    \item \textbf{Durabilité} : pour la baseline, chaque transaction exécutée est journalisée de manière synchrone dans un fichier CSV, sur le chemin critique. Cette approche naïve, génératrice de latence, sera remise en cause au Jalon 6 par une persistance \textit{Zero-Copy}.
\end{itemize}

\subsubsection{Modèle de données}

Les données fondamentales du \textit{matching engine}, sans optimisation, sont les suivantes :

\begin{itemize}
    \item \texttt{Side} (énumération)
        \begin{itemize}
            \item \texttt{BUY}, \texttt{SELL}
        \end{itemize}
    \item \texttt{OrderCommand} — représente une intention émise par le générateur d'ordres
        \begin{itemize}
            \item \texttt{commandType}: enum (\texttt{NEW}, \texttt{CANCEL}, \texttt{MODIFY})
            \item \texttt{orderId}: Long
            \item \texttt{side}: Side
            \item \texttt{price}: Double \textit{(borné, cf. exigences non-fonctionnelles)}
            \item \texttt{quantity}: Integer \textit{(quantité initiale pour \texttt{NEW}, nouvelle quantité pour \texttt{MODIFY}, ignorée pour \texttt{CANCEL})}
            \item \texttt{createdAt}: Long
        \end{itemize}
    \item \texttt{Order} — état d'un ordre tel que conservé dans le carnet
        \begin{itemize}
            \item \texttt{orderId}: Long
            \item \texttt{side}: Side
            \item \texttt{price}: Double
            \item \texttt{originalQuantity}: Integer
            \item \texttt{remainingQuantity}: Integer
            \item \texttt{createdAt}: Long
            \item \texttt{status}: enum (\texttt{RESTING}, \texttt{PARTIALLY\_FILLED}, \texttt{FILLED}, \texttt{CANCELLED})
        \end{itemize}
    \item \texttt{Trade} — résultat d'un \textit{match}
        \begin{itemize}
            \item \texttt{tradeId}: Long
            \item \texttt{buyOrderId}: Long
            \item \texttt{sellOrderId}: Long
            \item \texttt{executionPrice}: Double \textit{(prix de l'ordre résident, cf. règle de \textit{price-time priority})}
            \item \texttt{quantity}: Int
            \item \texttt{executedAt}: Long
        \end{itemize}
\end{itemize}

\textit{Source: Stock Exchange System Design | Order Matching, Trading Engine \& Scalability Explained - Youtube}

\subsubsection{Architecture du matching engine}

Notre objectif étant d'optimiser la latence de notre programme, nous minimisons les communications réseau et inter-processus. Nous optons pour une architecture monolithique modulaire à trois modules, communiquant via une \texttt{LinkedBlockingQueue} — une structure bloquante standard du JDK, choisie délibérément pour la baseline afin de fournir un point de comparaison avec le \textit{Ring Buffer} introduit au Jalon 1 :

\begin{itemize}
    \item \textbf{Module de génération} : produit des \texttt{OrderCommand} (\texttt{NEW}, \texttt{CANCEL}, \texttt{MODIFY}) à intervalles aléatoires et les dépose dans la file partagée à destination du module de \textit{matching}.
    \item \textbf{Module de \textit{matching}} : consomme les \texttt{OrderCommand}, maintient le carnet d'ordres et applique le \textit{price-time priority}. Le carnet est représenté par deux \texttt{TreeMap<Double, Deque<Order>>} (un par côté), structure standard du JDK offrant une recherche en $O(\log n)$ mais dont le coût en allocations et en \textit{pointer chasing} constituera la référence à améliorer dans les jalons suivants. L'accès concurrent au carnet est protégé par un verrou (\texttt{synchronized} ou \texttt{ReentrantLock}) — un choix délibérément coûteux, qui sera supprimé au profit du modèle \textit{Single-Writer} du Disruptor (Jalon 1).
    \item \textbf{Module de sortie et de persistance} : reçoit les \texttt{Trade} produits par le module de \textit{matching}, les affiche sur la sortie standard et les journalise de manière synchrone dans un fichier CSV.
\end{itemize}

\textit{Source: How Does a Matching Engine Work - Juan D Bustamante - Devnexus 2022 - Youtube}